\section{Ejercicios}

\subsection{Nivel básico}

\begin{ejercicio}
Dos cargas puntuales \(q_1=+\SI{2.0}{\micro\coulomb}\) y
\(q_2=+\SI{3.0}{\micro\coulomb}\) se encuentran separadas
\(\SI{0.50}{\metre}\).

Determine la magnitud de la fuerza eléctrica entre ellas e indique si la
interacción es atractiva o repulsiva.
\end{ejercicio}

\begin{respuesta}
\[
\boxed{F\approx\SI{0.216}{\newton}}
\]

La interacción es repulsiva.
\end{respuesta}

\begin{explicacion}
Como ambas cargas son positivas, tienen el mismo signo y se repelen.

Convertimos al SI:

\[
q_1=\SI{2.0e-6}{\coulomb}, \qquad q_2=\SI{3.0e-6}{\coulomb}.
\]

Aplicamos la ley de Coulomb:

\[
F=k\frac{|q_1q_2|}{r^2}.
\]

Sustituyendo,

\[
F=(8.99\times10^9)
\frac{(2.0\times10^{-6})(3.0\times10^{-6})}{(0.50)^2}
\approx\SI{0.216}{\newton}.
\]

Por tanto,

\[
\boxed{F\approx\SI{0.216}{\newton}}.
\]

La fuerza sobre cada carga actúa sobre la recta que las une y alejándose de
la otra carga.
\end{explicacion}


\begin{ejercicio}
Dos cargas puntuales \(q_1=+\SI{4.0}{\micro\coulomb}\) y
\(q_2=-\SI{6.0}{\micro\coulomb}\) están separadas
\(\SI{0.30}{\metre}\).

Calcule la magnitud de la fuerza eléctrica e indique su carácter.
\end{ejercicio}

\begin{respuesta}
\[
\boxed{F\approx\SI{2.40}{\newton}}
\]

La interacción es atractiva.
\end{respuesta}

\begin{explicacion}
Las cargas tienen signos opuestos, por lo que se atraen.

En unidades SI,

\[
q_1=\SI{4.0e-6}{\coulomb}, \qquad q_2=-\SI{6.0e-6}{\coulomb}.
\]

Aplicamos

\[
F=k\frac{|q_1q_2|}{r^2}.
\]

Entonces,

\[
F=(8.99\times10^9)
\frac{(4.0\times10^{-6})(6.0\times10^{-6})}{(0.30)^2}
\approx\SI{2.40}{\newton}.
\]

Por tanto,

\[
\boxed{F\approx\SI{2.40}{\newton}}.
\]

El signo de las cargas determina que la interacción es atractiva; la magnitud
de la fuerza es positiva.
\end{explicacion}


\begin{ejercicio}
Dos cargas puntuales de magnitudes
\[
|q_1|=\SI{2.0}{\micro\coulomb}, \qquad
|q_2|=\SI{8.0}{\micro\coulomb}
\]
ejercen entre sí una fuerza de magnitud \(\SI{0.80}{\newton}\).

Determine la distancia que las separa.
\end{ejercicio}

\begin{respuesta}
\[
\boxed{r\approx\SI{0.424}{\metre}}
\]
\end{respuesta}

\begin{explicacion}
Partimos de

\[
F=k\frac{|q_1q_2|}{r^2}.
\]

Despejamos:

\[
r=\sqrt{\frac{k|q_1q_2|}{F}}.
\]

Convertimos las cargas y sustituimos:

\[
r=
\sqrt{
\frac{(8.99\times10^9)(2.0\times10^{-6})(8.0\times10^{-6})}{0.80}
}.
\]

Por tanto,

\[
\boxed{r\approx\SI{0.424}{\metre}}.
\]
\end{explicacion}


\begin{ejercicio}
Una carga \(q_1=+\SI{5.0}{\micro\coulomb}\) se encuentra a
\(\SI{0.20}{\metre}\) de otra carga desconocida \(q_2\). La fuerza entre
ellas tiene magnitud \(\SI{2.25}{\newton}\) y es atractiva.

Determine \(q_2\).
\end{ejercicio}

\begin{respuesta}
\[
\boxed{q_2\approx-\SI{2.00}{\micro\coulomb}}
\]
\end{respuesta}

\begin{explicacion}
Como la interacción es atractiva y \(q_1\) es positiva, \(q_2\) debe ser
negativa.

De

\[
F=k\frac{|q_1q_2|}{r^2}
\]

despejamos

\[
|q_2|=\frac{Fr^2}{k|q_1|}.
\]

Sustituyendo,

\[
|q_2|=
\frac{(2.25)(0.20)^2}
{(8.99\times10^9)(5.0\times10^{-6})}
\approx\SI{2.00e-6}{\coulomb}.
\]

Como debe tener signo opuesto al de \(q_1\),

\[
\boxed{q_2\approx-\SI{2.00}{\micro\coulomb}}.
\]
\end{explicacion}


\begin{ejercicio}
Dos cargas ejercen inicialmente una fuerza de magnitud \(F\).

Si la distancia entre ellas se duplica sin modificar las cargas, ¿cuál será
la nueva fuerza?
\end{ejercicio}

\begin{respuesta}
\[
\boxed{F'=\frac{F}{4}}
\]
\end{respuesta}

\begin{explicacion}
La ley de Coulomb establece

\[
F\propto\frac{1}{r^2}.
\]

Si \(r'=2r\),

\[
\frac{F'}{F}=\frac{r^2}{(2r)^2}=\frac14.
\]

Por tanto,

\[
\boxed{F'=\frac{F}{4}}.
\]
\end{explicacion}


\begin{ejercicio}
Dos cargas ejercen inicialmente una fuerza de magnitud \(F\).

Si la primera carga se duplica y la segunda se triplica, manteniendo
constante la distancia, determine la nueva fuerza.
\end{ejercicio}

\begin{respuesta}
\[
\boxed{F'=6F}
\]
\end{respuesta}

\begin{explicacion}
La fuerza es proporcional al producto de las magnitudes de las cargas:

\[
F\propto|q_1q_2|.
\]

Si

\[
q_1'=2q_1, \qquad q_2'=3q_2,
\]

entonces

\[
|q_1'q_2'|=|(2q_1)(3q_2)|=6|q_1q_2|.
\]

Por tanto,

\[
\boxed{F'=6F}.
\]
\end{explicacion}


\begin{ejercicio}
La carga \(A\) ejerce sobre la carga \(B\) una fuerza eléctrica de magnitud
\(\SI{0.40}{\newton}\) hacia la derecha.

¿Cuál es la fuerza que \(B\) ejerce sobre \(A\)?
\end{ejercicio}

\begin{respuesta}
\[
\boxed{\SI{0.40}{\newton}\text{ hacia la izquierda}}
\]
\end{respuesta}

\begin{explicacion}
Por la tercera ley de Newton, las fuerzas que dos cargas ejercen mutuamente
tienen la misma magnitud y sentidos opuestos.

Si tomamos hacia la derecha como sentido positivo,

\[
F_{A\to B,x}=+\SI{0.40}{\newton}.
\]

Por tanto,

\[
F_{B\to A,x}=-\SI{0.40}{\newton}.
\]

Así, la fuerza que \(B\) ejerce sobre \(A\) tiene magnitud

\[
\boxed{\SI{0.40}{\newton}}
\]

y apunta hacia la izquierda.
\end{explicacion}


\begin{ejercicio}
Una carga positiva \(q_1\) se encuentra a la izquierda de una carga negativa
\(q_2\).

Indique el sentido de la fuerza eléctrica sobre cada carga.
\end{ejercicio}

\begin{respuesta}
La fuerza sobre \(q_1\) apunta hacia la derecha y la fuerza sobre \(q_2\)
hacia la izquierda.
\end{respuesta}

\begin{explicacion}
Las cargas tienen signos opuestos, por lo que se atraen.

La carga \(q_1\) es atraída hacia \(q_2\), por lo que su fuerza apunta hacia
la derecha.

La carga \(q_2\) es atraída hacia \(q_1\), por lo que su fuerza apunta hacia
la izquierda.

Ambas fuerzas tienen la misma magnitud.
\end{explicacion}


\begin{ejercicio}
Dos cargas \(q_1=\SI{25}{\nano\coulomb}\) y
\(q_2=\SI{40}{\nano\coulomb}\) se encuentran separadas
\(\SI{5.0}{\centi\metre}\).

Determine la magnitud de la fuerza eléctrica entre ellas.
\end{ejercicio}

\begin{respuesta}
\[
\boxed{F\approx\SI{3.60e-3}{\newton}}
\]
\end{respuesta}

\begin{explicacion}
Convertimos al SI:

\[
q_1=\SI{25e-9}{\coulomb}, \qquad
q_2=\SI{40e-9}{\coulomb}, \qquad
r=\SI{0.050}{\metre}.
\]

Aplicamos

\[
F=k\frac{|q_1q_2|}{r^2}.
\]

Sustituyendo,

\[
F=(8.99\times10^9)
\frac{(25\times10^{-9})(40\times10^{-9})}{(0.050)^2}.
\]

Por tanto,

\[
\boxed{F\approx\SI{3.60e-3}{\newton}}.
\]
\end{explicacion}


\begin{ejercicio}
Dos esferas cargadas tienen radios de aproximadamente
\(\SI{1}{\milli\metre}\) y sus centros están separados
\(\SI{0.50}{\metre}\).

¿Es razonable modelarlas como cargas puntuales al aplicar la ley de Coulomb?
\end{ejercicio}

\begin{respuesta}
Sí, la aproximación de carga puntual es razonable.
\end{respuesta}

\begin{explicacion}
El tamaño característico de cada esfera es del orden de

\[
L\sim\SI{1e-3}{\metre},
\]

mientras que

\[
r=\SI{0.50}{\metre}.
\]

Por tanto,

\[
L\ll r.
\]

Como las dimensiones de las esferas son muy pequeñas comparadas con la
separación entre ellas, es razonable tratarlas como cargas puntuales.
\end{explicacion}


\subsection{Nivel intermedio}

\begin{ejercicio}
Una carga \(q_0=+\SI{2.0}{\micro\coulomb}\) se encuentra en \(x=0\).
Otras dos cargas están ubicadas en

\[
q_1=+\SI{3.0}{\micro\coulomb}, \qquad x_1=-\SI{0.40}{\metre},
\]

y

\[
q_2=+\SI{8.0}{\micro\coulomb}, \qquad x_2=+\SI{0.80}{\metre}.
\]

Determine la fuerza neta sobre \(q_0\).
\end{ejercicio}

\begin{respuesta}
\[
\boxed{F_x\approx+\SI{0.112}{\newton}}
\]

La fuerza neta apunta hacia la derecha.
\end{respuesta}

\begin{explicacion}
La carga \(q_1\) repele a \(q_0\) hacia la derecha:

\[
F_1=k\frac{|q_1q_0|}{(0.40)^2}
\approx\SI{0.337}{\newton}.
\]

Por tanto,

\[
F_{1x}=+\SI{0.337}{\newton}.
\]

La carga \(q_2\) repele a \(q_0\) hacia la izquierda:

\[
F_2=k\frac{|q_2q_0|}{(0.80)^2}
\approx\SI{0.225}{\newton}.
\]

Así,

\[
F_{2x}=-\SI{0.225}{\newton}.
\]

Aplicamos superposición sobre el eje \(x\):

\[
F_x=F_{1x}+F_{2x}=0.337-0.225.
\]

Por tanto,

\[
\boxed{F_x\approx+\SI{0.112}{\newton}}.
\]

El signo positivo indica que la fuerza neta apunta hacia la derecha.
\end{explicacion}


\begin{ejercicio}
Una carga \(q_0=+\SI{1.0}{\micro\coulomb}\) se encuentra en el origen.
A su izquierda, en \(x=-\SI{0.30}{\metre}\), está
\(q_1=+\SI{4.0}{\micro\coulomb}\), y a su derecha, en
\(x=+\SI{0.60}{\metre}\), está \(q_2=+\SI{16}{\micro\coulomb}\).

Determine la fuerza neta sobre \(q_0\).
\end{ejercicio}

\begin{respuesta}
\[
\boxed{F_x=0}
\]
\end{respuesta}

\begin{explicacion}
La carga \(q_1\) repele a \(q_0\) hacia la derecha:

\[
F_1=k\frac{(4.0\times10^{-6})(1.0\times10^{-6})}{(0.30)^2}
\approx\SI{0.400}{\newton}.
\]

Por tanto,

\[
F_{1x}=+\SI{0.400}{\newton}.
\]

La carga \(q_2\) repele a \(q_0\) hacia la izquierda:

\[
F_2=k\frac{(16\times10^{-6})(1.0\times10^{-6})}{(0.60)^2}
\approx\SI{0.400}{\newton},
\]

de modo que

\[
F_{2x}=-\SI{0.400}{\newton}.
\]

Entonces,

\[
F_x=F_{1x}+F_{2x}=0.
\]

Por tanto,

\[
\boxed{F_x=0}.
\]

Las fuerzas existen, pero se cancelan.
\end{explicacion}


\begin{ejercicio}
Una carga \(q_0=+\SI{2.0}{\micro\coulomb}\) se encuentra en el origen.
Otra carga \(q_1=+\SI{3.0}{\micro\coulomb}\) está en
\((\SI{0.40}{\metre},0)\), y \(q_2=-\SI{4.0}{\micro\coulomb}\) está en
\((0,\SI{0.30}{\metre})\).

Determine las componentes rectangulares, la magnitud y la dirección de la
fuerza neta sobre \(q_0\).
\end{ejercicio}

\begin{respuesta}
\[
\boxed{F_x=-\SI{0.337}{\newton}}, \qquad
\boxed{F_y=+\SI{0.799}{\newton}}
\]

\[
\boxed{F_{\text{neta}}\approx\SI{0.867}{\newton}}
\]

\[
\boxed{\theta\approx112.9^\circ\text{ desde }+x}
\]
\end{respuesta}

\begin{explicacion}
La carga \(q_1\) es positiva y está a la derecha de \(q_0\), por lo que la
repele hacia la izquierda.

\[
F_1=k\frac{|q_1q_0|}{(0.40)^2}
\approx\SI{0.337}{\newton}.
\]

Sus componentes son

\[
F_{1x}=-\SI{0.337}{\newton}, \qquad F_{1y}=0.
\]

La carga \(q_2\) es negativa y está arriba de \(q_0\), por lo que atrae a
\(q_0\) hacia arriba.

\[
F_2=k\frac{|q_2q_0|}{(0.30)^2}
\approx\SI{0.799}{\newton}.
\]

Sus componentes son

\[
F_{2x}=0, \qquad F_{2y}=+\SI{0.799}{\newton}.
\]

Aplicamos superposición por componentes:

\[
F_x=F_{1x}+F_{2x}=-\SI{0.337}{\newton},
\]

\[
F_y=F_{1y}+F_{2y}=+\SI{0.799}{\newton}.
\]

La magnitud es

\[
F_{\text{neta}}=\sqrt{F_x^2+F_y^2}
=\sqrt{(-0.337)^2+(0.799)^2}
\approx\SI{0.867}{\newton}.
\]

Como \(F_x<0\) y \(F_y>0\), la fuerza se encuentra en el segundo cuadrante.

El ángulo de referencia es

\[
\theta_r=\tan^{-1}\left(\frac{|F_y|}{|F_x|}\right)
=\tan^{-1}\left(\frac{0.799}{0.337}\right)
\approx67.1^\circ.
\]

Por tanto,

\[
\boxed{\theta=180^\circ-67.1^\circ\approx112.9^\circ}.
\]
\end{explicacion}


\begin{ejercicio}
Sobre una carga actúan dos fuerzas eléctricas de magnitudes
\(F_1=\SI{0.50}{\newton}\) y \(F_2=\SI{0.80}{\newton}\), que forman entre sí
un ángulo de \(120^\circ\).

Determine la magnitud de la fuerza resultante utilizando componentes
rectangulares.
\end{ejercicio}

\begin{respuesta}
\[
\boxed{F_R=\SI{0.70}{\newton}}
\]
\end{respuesta}

\begin{explicacion}
Elegimos el eje \(x\) en la dirección de \(F_1\). Entonces,

\[
F_{1x}=\SI{0.50}{\newton}, \qquad F_{1y}=0.
\]

Para \(F_2\),

\[
F_{2x}=F_2\cos120^\circ=(0.80)(-0.5)
=-\SI{0.40}{\newton},
\]

\[
F_{2y}=F_2\sin120^\circ=(0.80)(0.866)
\approx+\SI{0.693}{\newton}.
\]

Sumamos las componentes:

\[
F_x=0.50-0.40=\SI{0.10}{\newton},
\]

\[
F_y=0+0.693=\SI{0.693}{\newton}.
\]

La magnitud de la resultante es

\[
F_R=\sqrt{F_x^2+F_y^2}
=\sqrt{(0.10)^2+(0.693)^2}.
\]

Por tanto,

\[
\boxed{F_R\approx\SI{0.70}{\newton}}.
\]
\end{explicacion}


\begin{ejercicio}
Una carga \(q_1=+\SI{2.0}{\micro\coulomb}\) se encuentra en el origen y otra
carga \(q_2=-\SI{5.0}{\micro\coulomb}\) se encuentra en
\((\SI{0.30}{\metre},\SI{0.40}{\metre})\).

Determine las componentes rectangulares y la magnitud de la fuerza que
\(q_1\) ejerce sobre \(q_2\).
\end{ejercicio}

\begin{respuesta}
\[
\boxed{F_x=-\SI{0.216}{\newton}}, \qquad
\boxed{F_y=-\SI{0.288}{\newton}}
\]

\[
\boxed{F\approx\SI{0.360}{\newton}}
\]
\end{respuesta}

\begin{explicacion}
La separación entre las cargas es

\[
r=\sqrt{(0.30)^2+(0.40)^2}=\SI{0.50}{\metre}.
\]

La magnitud de la fuerza es

\[
F=k\frac{|q_1q_2|}{r^2}
=(8.99\times10^9)
\frac{(2.0\times10^{-6})(5.0\times10^{-6})}{(0.50)^2}
\approx\SI{0.360}{\newton}.
\]

Como las cargas tienen signos opuestos, \(q_2\) es atraída hacia el origen.
Por tanto, ambas componentes son negativas.

De la geometría,

\[
\cos\theta=\frac{0.30}{0.50}=0.60, \qquad
\sin\theta=\frac{0.40}{0.50}=0.80.
\]

Así,

\[
F_x=-F\cos\theta=-(0.360)(0.60)
=-\SI{0.216}{\newton},
\]

\[
F_y=-F\sin\theta=-(0.360)(0.80)
=-\SI{0.288}{\newton}.
\]

Por tanto,

\[
\boxed{F_x=-\SI{0.216}{\newton}}, \qquad
\boxed{F_y=-\SI{0.288}{\newton}}.
\]
\end{explicacion}


\begin{ejercicio}
Una carga \(q_0=+\SI{2.0}{\micro\coulomb}\) se encuentra en el origen.
Dos cargas iguales \(q_1=q_2=+\SI{3.0}{\micro\coulomb}\) se encuentran en

\[
(-\SI{0.30}{\metre},+\SI{0.40}{\metre})
\]

y

\[
(-\SI{0.30}{\metre},-\SI{0.40}{\metre}).
\]

Determine la fuerza neta sobre \(q_0\).
\end{ejercicio}

\begin{respuesta}
\[
\boxed{F_x\approx+\SI{0.259}{\newton}}, \qquad
\boxed{F_y=0}
\]

La fuerza apunta hacia la derecha.
\end{respuesta}

\begin{explicacion}
Ambas cargas se encuentran a

\[
r=\sqrt{(0.30)^2+(0.40)^2}=\SI{0.50}{\metre}
\]

de \(q_0\).

Cada una ejerce una fuerza de magnitud

\[
F=k\frac{(3.0\times10^{-6})(2.0\times10^{-6})}{(0.50)^2}
\approx\SI{0.216}{\newton}.
\]

De la geometría,

\[
\cos\theta=\frac{0.30}{0.50}=0.60, \qquad
\sin\theta=\frac{0.40}{0.50}=0.80.
\]

Las componentes verticales tienen la misma magnitud y sentidos opuestos:

\[
F_y=F(0.80)-F(0.80)=0.
\]

Las componentes horizontales apuntan ambas hacia la derecha:

\[
F_x=2F(0.60)=2(0.216)(0.60)
\approx\SI{0.259}{\newton}.
\]

Por tanto,

\[
\boxed{F_x\approx+\SI{0.259}{\newton}}, \qquad
\boxed{F_y=0}.
\]
\end{explicacion}


\begin{ejercicio}
Dos cargas positivas \(q_1=+\SI{4.0}{\micro\coulomb}\) y
\(q_2=+\SI{9.0}{\micro\coulomb}\) se encuentran sobre el eje \(x\),
separadas \(\SI{1.0}{\metre}\).

Determine el punto entre ellas donde una tercera carga puntual no
experimentaría fuerza eléctrica neta.
\end{ejercicio}

\begin{respuesta}
\[
\boxed{x=\SI{0.40}{\metre}}
\]

medidos desde \(q_1\) hacia \(q_2\).
\end{respuesta}

\begin{explicacion}
Sea \(x\) la distancia desde \(q_1\).

Entre las cargas, las fuerzas sobre una tercera carga tienen sentidos
opuestos, por lo que pueden cancelarse.

La distancia hasta \(q_2\) es \(1-x\). Para el equilibrio,

\[
F_{1x}+F_{2x}=0.
\]

En magnitud,

\[
k\frac{|q_1q_0|}{x^2}
=
k\frac{|q_2q_0|}{(1-x)^2}.
\]

Cancelando \(k\) y \(|q_0|\),

\[
\frac{4}{x^2}=\frac{9}{(1-x)^2}.
\]

Tomando raíces positivas,

\[
\frac{2}{x}=\frac{3}{1-x}.
\]

Entonces,

\[
2(1-x)=3x
\]

y

\[
\boxed{x=\SI{0.40}{\metre}}.
\]

El punto se encuentra más cerca de la carga de menor magnitud.
\end{explicacion}


\begin{ejercicio}
Dos cargas \(q_1=+\SI{4.0}{\micro\coulomb}\) y
\(q_2=-\SI{9.0}{\micro\coulomb}\) están ubicadas en \(x=0\) y
\(x=\SI{1.0}{\metre}\), respectivamente.

Determine sobre el eje \(x\) un punto donde una tercera carga pueda
experimentar fuerza eléctrica neta nula.
\end{ejercicio}

\begin{respuesta}
\[
\boxed{x=-\SI{2.0}{\metre}}
\]
\end{respuesta}

\begin{explicacion}
Entre las cargas, las dos fuerzas tienen el mismo sentido, por lo que no
pueden cancelarse.

El punto debe encontrarse fuera del intervalo, del lado de la carga de menor
magnitud.

Sea \(a\) la distancia a la izquierda de \(q_1\). Entonces,

\[
r_1=a, \qquad r_2=1+a.
\]

Para que la componente total sea nula,

\[
F_{1x}+F_{2x}=0.
\]

Por tanto, las magnitudes deben ser iguales:

\[
k\frac{4|q_0|}{a^2}
=
k\frac{9|q_0|}{(1+a)^2}.
\]

Simplificando,

\[
\frac{2}{a}=\frac{3}{1+a}.
\]

Entonces,

\[
2(1+a)=3a,
\]

de donde

\[
a=2.
\]

Como el punto está a la izquierda del origen,

\[
\boxed{x=-\SI{2.0}{\metre}}.
\]
\end{explicacion}


\begin{ejercicio}
Una carga \(q_0=+\SI{1.0}{\micro\coulomb}\) se encuentra en el origen.
Una carga \(q_1=+\SI{4.0}{\micro\coulomb}\) se encuentra en
\(x=-\SI{0.20}{\metre}\).

¿Qué carga positiva \(q_2\), ubicada en \(x=+\SI{0.40}{\metre}\), hará que
la fuerza neta sobre \(q_0\) sea \(\SI{0.4495}{\newton}\) hacia la derecha?
\end{ejercicio}

\begin{respuesta}
\[
\boxed{q_2=\SI{8.0}{\micro\coulomb}}
\]
\end{respuesta}

\begin{explicacion}
La carga \(q_1\) repele a \(q_0\) hacia la derecha:

\[
F_{1x}=+\SI{0.899}{\newton}.
\]

La carga \(q_2\), al ser positiva y estar a la derecha, repele a \(q_0\)
hacia la izquierda, por lo que

\[
F_{2x}<0.
\]

Aplicamos la suma de componentes:

\[
F_x=F_{1x}+F_{2x}.
\]

En magnitudes,

\[
0.4495=0.899-F_2,
\]

de donde

\[
F_2=\SI{0.4495}{\newton}.
\]

Ahora,

\[
F_2=k\frac{|q_2q_0|}{(0.40)^2}.
\]

Despejando,

\[
q_2=
\frac{F_2(0.40)^2}{kq_0}
=
\frac{(0.4495)(0.40)^2}
{(8.99\times10^9)(1.0\times10^{-6})}.
\]

Finalmente,

\[
\boxed{q_2=\SI{8.0}{\micro\coulomb}}.
\]
\end{explicacion}


\begin{ejercicio}
Dos cargas ejercen inicialmente una fuerza de magnitud
\(\SI{1.20}{\newton}\).

Después, la primera carga se reduce a la mitad, la segunda se triplica y la
distancia aumenta hasta \(1.5\) veces su valor inicial.

Determine la nueva fuerza.
\end{ejercicio}

\begin{respuesta}
\[
\boxed{F'=\SI{0.80}{\newton}}
\]
\end{respuesta}

\begin{explicacion}
Utilizamos

\[
\frac{F'}{F}
=
\frac{|q_1'q_2'|}{|q_1q_2|}
\frac{r^2}{r'^2}.
\]

Como

\[
q_1'=\frac12q_1, \qquad q_2'=3q_2, \qquad r'=1.5r,
\]

se obtiene

\[
\frac{F'}{F}
=
\left(\frac12\right)(3)\frac{1}{(1.5)^2}
=
\frac23.
\]

Por tanto,

\[
F'=\frac23(1.20)=\SI{0.80}{\newton}.
\]

\[
\boxed{F'=\SI{0.80}{\newton}}.
\]
\end{explicacion}


\subsection{Nivel avanzado}

\begin{ejercicio}
Cuatro cargas iguales \(q=+\SI{1.0}{\micro\coulomb}\) se encuentran en los
cuatro vértices de un cuadrado de lado \(a=\SI{0.20}{\metre}\).

Determine las componentes rectangulares, la magnitud y la dirección de la
fuerza eléctrica neta sobre la carga situada en el vértice inferior
izquierdo.
\end{ejercicio}

\begin{respuesta}
\[
\boxed{F_x\approx-\SI{0.304}{\newton}}, \qquad
\boxed{F_y\approx-\SI{0.304}{\newton}}
\]

\[
\boxed{F_{\text{neta}}\approx\SI{0.430}{\newton}}
\]

La fuerza apunta a \(45^\circ\) por debajo del eje \(-x\).
\end{respuesta}

\begin{explicacion}
Ubicamos la carga de interés en el vértice inferior izquierdo.

Las dos cargas adyacentes se encuentran a distancia \(a\). Cada una produce

\[
F_a=k\frac{q^2}{a^2}
=(8.99\times10^9)
\frac{(1.0\times10^{-6})^2}{(0.20)^2}
\approx\SI{0.225}{\newton}.
\]

La carga situada a la derecha produce

\[
F_{ax}=-\SI{0.225}{\newton}, \qquad F_{ay}=0,
\]

mientras que la carga situada arriba produce

\[
F_{bx}=0, \qquad F_{by}=-\SI{0.225}{\newton}.
\]

La carga diagonal está a

\[
r=a\sqrt2.
\]

Su fuerza tiene magnitud

\[
F_d=k\frac{q^2}{2a^2}\approx\SI{0.112}{\newton}.
\]

Como forma \(45^\circ\) con ambos ejes y también es repulsiva,

\[
F_{dx}=-F_d\cos45^\circ\approx-\SI{0.0795}{\newton},
\]

\[
F_{dy}=-F_d\sin45^\circ\approx-\SI{0.0795}{\newton}.
\]

Sumamos:

\[
F_x=-0.225-0.0795\approx-\SI{0.304}{\newton},
\]

\[
F_y=-0.225-0.0795\approx-\SI{0.304}{\newton}.
\]

La magnitud es

\[
F_{\text{neta}}
=\sqrt{F_x^2+F_y^2}
\approx\SI{0.430}{\newton}.
\]

Como las dos componentes son negativas e iguales, la fuerza se encuentra
sobre la diagonal del tercer cuadrante:

\[
\boxed{\theta=45^\circ\text{ por debajo del eje }-x}.
\]
\end{explicacion}


\begin{ejercicio}
Tres cargas iguales \(q=+\SI{2.0}{\micro\coulomb}\) se encuentran en los
vértices de un triángulo equilátero de lado \(a=\SI{0.30}{\metre}\).

Determine la fuerza neta sobre la carga situada en el vértice superior.
\end{ejercicio}

\begin{respuesta}
\[
\boxed{F_x=0}, \qquad
\boxed{F_y\approx+\SI{0.692}{\newton}}
\]

\[
\boxed{F_{\text{neta}}\approx\SI{0.692}{\newton}}
\]

La fuerza apunta verticalmente hacia arriba.
\end{respuesta}

\begin{explicacion}
Cada carga inferior repele a la carga superior con magnitud

\[
F=k\frac{q^2}{a^2}
=(8.99\times10^9)
\frac{(2.0\times10^{-6})^2}{(0.30)^2}
\approx\SI{0.400}{\newton}.
\]

Cada fuerza forma \(60^\circ\) con la horizontal.

Las componentes horizontales son

\[
F_{1x}=-F\cos60^\circ, \qquad
F_{2x}=+F\cos60^\circ,
\]

por lo que

\[
F_x=0.
\]

Las componentes verticales se suman:

\[
F_y=2F\sin60^\circ.
\]

Entonces,

\[
F_y=2(0.400)(0.866)\approx\SI{0.692}{\newton}.
\]

Por tanto,

\[
\boxed{F_x=0}, \qquad
\boxed{F_y\approx+\SI{0.692}{\newton}}.
\]
\end{explicacion}


\begin{ejercicio}
Una carga \(q_0=-\SI{2.0}{\micro\coulomb}\) se encuentra en
\((1,1)\,\si{\metre}\).

Una carga \(q_1=+\SI{3.0}{\micro\coulomb}\) está en
\((-2,1)\,\si{\metre}\), y otra \(q_2=-\SI{4.0}{\micro\coulomb}\) está en
\((1,-2)\,\si{\metre}\).

Determine las componentes rectangulares, la magnitud y la dirección de la
fuerza neta sobre \(q_0\).
\end{ejercicio}

\begin{respuesta}
\[
\boxed{F_x=-\SI{5.99e-3}{\newton}}, \qquad
\boxed{F_y=+\SI{7.99e-3}{\newton}}
\]

\[
\boxed{F_{\text{neta}}\approx\SI{9.99e-3}{\newton}}
\]

\[
\boxed{\theta\approx126.9^\circ\text{ desde }+x}
\]
\end{respuesta}

\begin{explicacion}
La distancia entre \(q_0\) y \(q_1\) es

\[
r_1=\SI{3}{\metre}.
\]

Como \(q_0\) es negativa y \(q_1\) positiva, la interacción es atractiva y
la fuerza apunta hacia la izquierda:

\[
F_1=k\frac{(2.0\times10^{-6})(3.0\times10^{-6})}{3^2}
\approx\SI{5.99e-3}{\newton}.
\]

Por tanto,

\[
F_{1x}=-\SI{5.99e-3}{\newton}, \qquad F_{1y}=0.
\]

La distancia entre \(q_0\) y \(q_2\) también es

\[
r_2=\SI{3}{\metre}.
\]

Como ambas cargas son negativas, se repelen. Como \(q_2\) está debajo de
\(q_0\), la fuerza apunta hacia arriba:

\[
F_2=k\frac{(2.0\times10^{-6})(4.0\times10^{-6})}{3^2}
\approx\SI{7.99e-3}{\newton}.
\]

Entonces,

\[
F_{2x}=0, \qquad F_{2y}=+\SI{7.99e-3}{\newton}.
\]

Sumamos por componentes:

\[
F_x=-\SI{5.99e-3}{\newton},
\qquad
F_y=+\SI{7.99e-3}{\newton}.
\]

La magnitud es

\[
F_{\text{neta}}
=
\sqrt{F_x^2+F_y^2}
\approx\SI{9.99e-3}{\newton}.
\]

La fuerza se encuentra en el segundo cuadrante. El ángulo de referencia es

\[
\theta_r=\tan^{-1}\left(\frac{7.99}{5.99}\right)\approx53.1^\circ.
\]

Por tanto,

\[
\boxed{\theta\approx180^\circ-53.1^\circ=126.9^\circ}.
\]
\end{explicacion}


\begin{ejercicio}
Dos cargas positivas \(q_1\) y \(q_2\) están separadas una distancia \(L\).

Obtenga una expresión para la posición, medida desde \(q_1\), donde una
tercera carga experimenta fuerza neta nula.
\end{ejercicio}

\begin{respuesta}
\[
\boxed{x=\frac{L\sqrt{q_1}}{\sqrt{q_1}+\sqrt{q_2}}}
\]
\end{respuesta}

\begin{explicacion}
Como las cargas tienen el mismo signo, el punto de cancelación se encuentra
entre ellas.

Sea \(x\) la distancia desde \(q_1\). La distancia a \(q_2\) es \(L-x\).

Sobre el eje \(x\), la condición de equilibrio es

\[
F_{1x}+F_{2x}=0.
\]

Por tanto, sus magnitudes deben ser iguales:

\[
k\frac{|q_1q_0|}{x^2}
=
k\frac{|q_2q_0|}{(L-x)^2}.
\]

Cancelando \(k\) y \(|q_0|\),

\[
\frac{q_1}{x^2}=\frac{q_2}{(L-x)^2}.
\]

Tomando raíces,

\[
\frac{\sqrt{q_1}}{x}
=
\frac{\sqrt{q_2}}{L-x}.
\]

Entonces,

\[
\sqrt{q_1}(L-x)=\sqrt{q_2}\,x.
\]

Despejando,

\[
\boxed{x=\frac{L\sqrt{q_1}}{\sqrt{q_1}+\sqrt{q_2}}}.
\]
\end{explicacion}


\begin{ejercicio}
Una carga \(+q\) se encuentra en \(x=0\) y una carga \(-4q\) en \(x=L\).

Determine dónde, sobre el eje \(x\), puede colocarse una tercera carga para
que la fuerza eléctrica neta sobre ella sea cero.
\end{ejercicio}

\begin{respuesta}
\[
\boxed{x=-L}
\]
\end{respuesta}

\begin{explicacion}
Entre las cargas, las fuerzas sobre una tercera carga tienen el mismo
sentido, por lo que no pueden cancelarse.

El punto debe estar fuera del segmento y del lado de la carga de menor
magnitud.

Sea \(x_d\) la distancia a la izquierda de \(+q\). Entonces,

\[
r_1=x_d, \qquad r_2=L+x_d.
\]

La condición sobre el eje \(x\) es

\[
F_{1x}+F_{2x}=0.
\]

Por tanto,

\[
k\frac{|qq_0|}{x_d^2}
=
k\frac{|(-4q)q_0|}{(L+x_d)^2}.
\]

Simplificando,

\[
\frac{1}{x_d^2}
=
\frac{4}{(L+x_d)^2}.
\]

Tomando raíces,

\[
\frac{1}{x_d}=\frac{2}{L+x_d}.
\]

Entonces,

\[
L+x_d=2x_d,
\]

por lo que

\[
x_d=L.
\]

Como el punto está a la izquierda del origen,

\[
\boxed{x=-L}.
\]
\end{explicacion}


\begin{ejercicio}
Una carga \(q_0=+\SI{1.0}{\micro\coulomb}\) se encuentra en el origen.

Dos cargas \(q_1=q_2=+\SI{2.0}{\micro\coulomb}\) están ubicadas en

\[
(\SI{0.30}{\metre},+\SI{0.40}{\metre})
\]

y

\[
(\SI{0.30}{\metre},-\SI{0.40}{\metre}).
\]

Se coloca una tercera carga \(q_3\) en
\((-\SI{0.20}{\metre},0)\).

Determine el signo y la magnitud de \(q_3\) para que la fuerza neta sobre
\(q_0\) sea cero.
\end{ejercicio}

\begin{respuesta}
\[
\boxed{q_3=+\SI{0.384}{\micro\coulomb}}
\]
\end{respuesta}

\begin{explicacion}
La distancia desde \(q_1\) y \(q_2\) hasta \(q_0\) es

\[
r=\sqrt{(0.30)^2+(0.40)^2}=\SI{0.50}{\metre}.
\]

Cada fuerza tiene magnitud

\[
F=k\frac{(2.0\times10^{-6})(1.0\times10^{-6})}{(0.50)^2}
=\SI{0.07192}{\newton}.
\]

Por simetría, sus componentes verticales se cancelan:

\[
F_y=0.
\]

De la geometría,

\[
\cos\theta=\frac{0.30}{0.50}=0.60.
\]

Ambas componentes horizontales apuntan hacia la izquierda:

\[
F_{12,x}=-2F\cos\theta.
\]

Entonces,

\[
F_{12,x}=-2(0.07192)(0.60)
=-\SI{0.086304}{\newton}.
\]

Para que

\[
F_x=0,
\]

la carga \(q_3\) debe producir una componente

\[
F_{3x}=+\SI{0.086304}{\newton}.
\]

Como \(q_3\) está a la izquierda de \(q_0\), debe ser positiva para repelerla
hacia la derecha.

Por tanto,

\[
F_3=k\frac{q_3q_0}{(0.20)^2}.
\]

Despejando,

\[
q_3=
\frac{(0.086304)(0.20)^2}
{(8.99\times10^9)(1.0\times10^{-6})}
=3.84\times10^{-7}\,\si{\coulomb}.
\]

Así,

\[
\boxed{q_3=+\SI{0.384}{\micro\coulomb}}.
\]
\end{explicacion}


\begin{ejercicio}
Dos cargas \(q_1=+\SI{2.0}{\micro\coulomb}\) y
\(q_2=-\SI{3.0}{\micro\coulomb}\) están separadas
\(\SI{0.40}{\metre}\).

Un experimento afirma que la fuerza de \(q_1\) sobre \(q_2\) es
\(\SI{0.34}{\newton}\) hacia \(q_1\), mientras que la fuerza de \(q_2\)
sobre \(q_1\) es \(\SI{0.31}{\newton}\) hacia \(q_2\).

¿Son físicamente compatibles estos resultados si se consideran exactos?
\end{ejercicio}

\begin{respuesta}
No.

Ambas fuerzas deberían tener una magnitud aproximada de

\[
\boxed{\SI{0.337}{\newton}}
\]

y sentidos opuestos.
\end{respuesta}

\begin{explicacion}
La ley de Coulomb predice

\[
F=k\frac{|q_1q_2|}{r^2}.
\]

Sustituyendo,

\[
F=(8.99\times10^9)
\frac{(2.0\times10^{-6})(3.0\times10^{-6})}{(0.40)^2}
\approx\SI{0.337}{\newton}.
\]

Si elegimos el eje \(x\) sobre la línea que une las cargas, las dos fuerzas
deben cumplir

\[
F_{1x}=-F_{2x}.
\]

Es decir, sus componentes tienen igual magnitud y signos opuestos.

Los valores \(\SI{0.34}{\newton}\) y \(\SI{0.31}{\newton}\), considerados
exactos, no cumplen esta condición.

Si fueran mediciones experimentales, sería necesario conocer sus
incertidumbres antes de decidir si son compatibles.
\end{explicacion}


\begin{ejercicio}
Tres cargas ocupan los vértices de un triángulo equilátero de lado
\(a=\SI{0.30}{\metre}\).

En el vértice superior se encuentra \(q_0=+\SI{2.0}{\micro\coulomb}\), en el
inferior izquierdo \(q_1=+\SI{3.0}{\micro\coulomb}\) y en el inferior
derecho \(q_2=-\SI{6.0}{\micro\coulomb}\).

Determine las componentes rectangulares, la magnitud y la dirección de la
fuerza neta sobre \(q_0\).
\end{ejercicio}

\begin{respuesta}
\[
\boxed{F_x\approx+\SI{0.899}{\newton}}, \qquad
\boxed{F_y\approx-\SI{0.519}{\newton}}
\]

\[
\boxed{F_{\text{neta}}\approx\SI{1.04}{\newton}}
\]

\[
\boxed{\theta=30^\circ\text{ por debajo del eje }+x}
\]
\end{respuesta}

\begin{explicacion}
La fuerza producida por \(q_1\) es repulsiva:

\[
F_1=k\frac{|q_1q_0|}{a^2}
\approx\SI{0.599}{\newton}.
\]

Forma \(60^\circ\) con el eje \(+x\), por lo que

\[
F_{1x}=F_1\cos60^\circ, \qquad
F_{1y}=F_1\sin60^\circ.
\]

La fuerza producida por \(q_2\) es atractiva:

\[
F_2=k\frac{|q_2q_0|}{a^2}
\approx\SI{1.199}{\newton}.
\]

Apunta \(60^\circ\) por debajo del eje \(+x\):

\[
F_{2x}=F_2\cos60^\circ, \qquad
F_{2y}=-F_2\sin60^\circ.
\]

Sumando componentes,

\[
F_x=(0.599)(0.5)+(1.199)(0.5)
\approx\SI{0.899}{\newton},
\]

\[
F_y=(0.599)(0.866)-(1.199)(0.866)
\approx-\SI{0.519}{\newton}.
\]

La magnitud es

\[
F_{\text{neta}}
=
\sqrt{F_x^2+F_y^2}
\approx\SI{1.04}{\newton}.
\]

Como \(F_x>0\) y \(F_y<0\), la fuerza está en el cuarto cuadrante.

\[
\theta_r=
\tan^{-1}\left(\frac{|F_y|}{|F_x|}\right)
=
\tan^{-1}\left(\frac{0.519}{0.899}\right)
\approx30^\circ.
\]

Por tanto,

\[
\boxed{\theta=30^\circ\text{ por debajo del eje }+x}.
\]
\end{explicacion}


\begin{ejercicio}
Una carga \(q_0=+\SI{1.0}{\micro\coulomb}\) se encuentra en el origen.
Una carga \(q_1=+\SI{4.0}{\micro\coulomb}\) se encuentra en
\(x=-\SI{0.30}{\metre}\).

Otra carga \(q_2\), de signo desconocido, se encuentra en
\(x=+\SI{0.60}{\metre}\).

Se sabe que la fuerza neta sobre \(q_0\) es
\(\SI{0.200}{\newton}\) hacia la derecha.

Determine el signo y la magnitud de \(q_2\).
\end{ejercicio}

\begin{respuesta}
\[
\boxed{q_2\approx+\SI{7.99}{\micro\coulomb}}
\]
\end{respuesta}

\begin{explicacion}
La fuerza de \(q_1\) sobre \(q_0\) apunta hacia la derecha:

\[
F_{1x}=+
k\frac{|q_1q_0|}{(0.30)^2}
\approx+\SI{0.3996}{\newton}.
\]

La fuerza neta es

\[
F_x=+\SI{0.200}{\newton}.
\]

Por tanto, la fuerza producida por \(q_2\) debe apuntar hacia la izquierda:

\[
F_{2x}=F_x-F_{1x}
=0.200-0.3996
=-\SI{0.1996}{\newton}.
\]

Como \(q_2\) está a la derecha de \(q_0\), debe ser positiva para repelerla
hacia la izquierda.

Su magnitud satisface

\[
F_2=k\frac{|q_2q_0|}{(0.60)^2}.
\]

Despejando,

\[
q_2=
\frac{(0.1996)(0.60)^2}
{(8.99\times10^9)(1.0\times10^{-6})}.
\]

Por tanto,

\[
\boxed{q_2\approx+\SI{7.99}{\micro\coulomb}}.
\]
\end{explicacion}


\begin{ejercicio}
Una carga \(q_1=+\SI{2.0}{\micro\coulomb}\) se encuentra en el punto
\((1,-2,2)\,\si{\metre}\), y una carga
\(q_2=-\SI{3.0}{\micro\coulomb}\) se encuentra en
\((4,2,2)\,\si{\metre}\).

Determine las componentes rectangulares y la magnitud de la fuerza que
\(q_1\) ejerce sobre \(q_2\).
\end{ejercicio}

\begin{respuesta}
\[
\boxed{F_x=-\SI{1.29e-3}{\newton}}
\]

\[
\boxed{F_y=-\SI{1.73e-3}{\newton}}
\]

\[
\boxed{F_z=0}
\]

\[
\boxed{F\approx\SI{2.16e-3}{\newton}}
\]
\end{respuesta}

\begin{explicacion}
Las diferencias de coordenadas entre las cargas son

\[
\Delta x=4-1=\SI{3}{\metre},
\]

\[
\Delta y=2-(-2)=\SI{4}{\metre},
\]

\[
\Delta z=2-2=0.
\]

La distancia entre ellas es

\[
r=\sqrt{(\Delta x)^2+(\Delta y)^2+(\Delta z)^2}
=\sqrt{3^2+4^2}
=\SI{5}{\metre}.
\]

La magnitud de la fuerza es

\[
F=k\frac{|q_1q_2|}{r^2}
=(8.99\times10^9)
\frac{(2.0\times10^{-6})(3.0\times10^{-6})}{25}
\approx\SI{2.16e-3}{\newton}.
\]

Como las cargas tienen signos opuestos, la fuerza sobre \(q_2\) apunta desde
\(q_2\) hacia \(q_1\). Por ello, sus componentes \(x\) e \(y\) son negativas.

De la geometría,

\[
\frac{|\Delta x|}{r}=\frac35, \qquad
\frac{|\Delta y|}{r}=\frac45.
\]

Entonces,

\[
F_x=-F\frac{3}{5}
=-(2.16\times10^{-3})\frac35
\approx-\SI{1.29e-3}{\newton},
\]

\[
F_y=-F\frac{4}{5}
=-(2.16\times10^{-3})\frac45
\approx-\SI{1.73e-3}{\newton},
\]

y como no existe separación en \(z\),

\[
F_z=0.
\]

Por tanto,

\[
\boxed{F_x=-\SI{1.29e-3}{\newton}}, \qquad
\boxed{F_y=-\SI{1.73e-3}{\newton}}, \qquad
\boxed{F_z=0}.
\]

La magnitud es

\[
\boxed{F\approx\SI{2.16e-3}{\newton}}.
\]
\end{explicacion}